\chapter{曲线积分和曲面积分}
\begin{definition}{第一类、第二类曲线积分和曲面积分}{}
    设曲线分片$\cont^1$且正则，曲面分片光滑且具有不重复覆盖的正则参数化，以下被积函数连续。对平面曲线$\bm r(t)=(x(t),y(t))$，$\alpha\leq t\leq\beta$，第一类曲线积分为
    \[
        \int_Lf\d\ell=\int_\alpha^\beta f(x(t),y(t))\sqrt{x'(t)^2+y'(t)^2}\d t.
    \]
    空间曲线则用$\norm{\bm r'(t)}$作为弧长因子。第一类曲面积分为
    \[
        \iint_Sf\d\sigma=\iint_Df(\bm r(u,v))\norm{\bm r_u\times\bm r_v}\d u\nd v.
    \]
    第一类积分不依赖曲线或曲面的取向。第二类曲线积分沿指定路径$L$从$A$指向$B$；若$\bm F=(X,Y)$，则
    \[
        \int_L\bm F\cdot\d\bm\ell
        =\int_LX\d x+Y\d y
        =\int_\alpha^\beta\bm F(\bm r(t))\cdot\bm r'(t)\d t.
    \]
    积分通常依赖路径，而不仅由端点决定；反向后变号。若曲面可定向，选取连续的单位法向$\uvec{n}$，令$\bm V=(X,Y,Z)$，第二类曲面积分为
    \[
        \iint_{S^+}\bm V\cdot\d\bm S
        =\iint_{S^+}X\d y\wedge\d z+Y\d z\wedge\d x+Z\d x\wedge\d y.
    \]
\end{definition}
第二类曲面积分可写为第一类曲面积分
\[
    \iint_{S^+}\bm V\cdot\d\bm S=\iint_S\bm V\cdot\uvec{n}\d\sigma.
\]
记$\bm r_u\times\bm r_v=(A,B,C)$，计算公式为
\[
    \iint_{S^+}\bm V\cdot\d\bm S
    =\pm\iint_D\bigb{XA+YB+ZC}\d u\nd v.
\]
右侧各量均在参数点对应的曲面点取值；当$\bm r_u\times\bm r_v$与指定法向同向时取正号，反向时取负号。多块参数曲面应分别匹配取向后相加。
\section{Green, Gauss, Stokes公式}
取向约定可参见\href{https://openstax.org/books/calculus-volume-3/pages/6-7-stokes-theorem}{OpenStax Calculus Volume 3, §6.7}。
\begin{theorem}{Green公式}{}
    设$D$为有界平面区域，边界由有限条分片光滑的简单闭曲线组成；$X,Y$在$\overline D$的开邻域内属于$\cont^1$。边界采用正向：沿边界行进时区域在左侧，即外边界逆时针、内孔边界顺时针。则
    \begin{align*}
        \iint_D\Bigp{\pv Xx+\pv Yy}\d x\nd y&=\oint_{\p D}X\d y-Y\d x,\\
        \iint_D\Bigp{\pv Yx-\pv Xy}\d x\nd y&=\oint_{\p D}X\d x+Y\d y.
    \end{align*}
    区域可以有孔，但所有边界分支都必须计入。
\end{theorem}
\begin{theorem}{Gauss公式}{}
    设$\Omega\subset\RR^3$为有界区域，边界分片光滑，$\bm V=(X,Y,Z)$在$\overline\Omega$的开邻域内属于$\cont^1$。边界取相对于$\Omega$的外法向，包括空腔内壁，则
    \[
        \iiint_\Omega\div\bm V\d x\nd y\nd z
        =\iint_{\p\Omega}\bm V\cdot\uvec{n}\d\sigma.
    \]
\end{theorem}
\begin{theorem}{Stokes公式}{}
    设$S$为紧的、可定向的分片光滑曲面，边界$\p S$分片光滑，$\bm V$在$S$的开邻域内属于$\cont^1$。曲面法向与边界方向按右手规则相容，则
    \[
        \iint_S(\curl\bm V)\cdot\uvec{n}\d\sigma
        =\oint_{\p S}\bm V\cdot\d\bm\ell.
    \]
    翻转曲面法向时，边界方向也必须同时翻转。
\end{theorem}
